\section{Praktikum}
\label{sec:praktikum-bab14}

\begin{praktikumbox}
\textbf{Judul:} Operasi masukan/keluaran dengan \code{INT 21h} dan
penelusuran stack pada pemanggilan \asm{INT}.\\[0.3em]
\textbf{Sub-CPMK:} 2.4 \quad\textbullet\quad
\textbf{Alokasi waktu:} 100 menit \quad\textbullet\quad
\textbf{Bentuk:} individu \quad\textbullet\quad
\textbf{Tugas 7 dikumpulkan pada praktikum ini.}
\end{praktikumbox}

\subsection{Tujuan Praktikum}

Setelah menyelesaikan praktikum ini, mahasiswa mampu menulis program
masukan/keluaran dengan \code{INT 21h} dan menunjukkan isi stack
sebelum serta sesudah instruksi \asm{INT} dijalankan.

\subsection{Alat dan Bahan}

\begin{itemize}
    \item Komputer dengan Emu8086, GUI Turbo Assembler, atau DOSBox
          beserta DEBUG.
    \item Berkas \file{input-nama.asm} dan \file{warna-layar.asm} dari
          direktori \file{code/asm/}.
    \item Jurnal praktikum untuk mencatat nilai register dan isi
          stack.
\end{itemize}

\subsection{Langkah Kerja}

\subsubsection{A. Masukan/keluaran dengan \code{INT 21h}}

\begin{enumerate}
    \item Rakit dan jalankan \file{input-nama.asm} seperti pada
          Kode~\ref{lst:input-nama}, lalu catat keluarannya.
    \item Ubah kapasitas penyangga dari \dec{20} menjadi \dec{5}, lalu
          ketik nama yang lebih panjang dari empat karakter dan
          jelaskan hasilnya dengan mengacu pada
          Kode~\ref{lst:buffer-0a}.
    \item Hapus baris \lstasm|MOV BYTE PTR [SI], '$'|, rakit ulang, dan
          jalankan. Amati keluarannya dan jelaskan penyebabnya.
    \item Tambahkan satu baris keluaran yang mencetak NIM mahasiswa
          memakai fungsi \hex{09}.
\end{enumerate}

\subsubsection{B. Menelusuri Stack pada Pemanggilan \asm{INT}}

\begin{enumerate}
    \item Buka DOSBox, lakukan \file{mount} ke direktori kerja, lalu
          rakit program dengan \file{tasm} dan \file{tlink}.
    \item Jalankan \file{debug input-nama.exe} dan tampilkan instruksi
          program dengan perintah \file{U}. Tandai alamat instruksi
          \lstasm|INT 21h| yang memanggil fungsi \hex{09}.
    \item Tampilkan register dengan perintah \file{R}, lalu catat
          nilai \reg{SP}, \reg{CS}, \reg{IP}, dan \reg{FLAGS}.
    \item Jalankan \file{T} sampai tepat sebelum \lstasm|INT 21h|,
          catat nilai \reg{SP}, lalu tampilkan puncak stack dengan
          \file{D SS:SP L 6} dan catat enam byte yang tampak.
    \item Jalankan \file{T} satu kali untuk mengeksekusi
          \lstasm|INT 21h|, lalu ulangi perintah \file{D SS:SP L 6}.
    \item Bandingkan hasil langkah 4 dan 5, jelaskan mengapa \reg{SP}
          berubah, dan tentukan nilai mana yang merupakan \reg{IP},
          \reg{CS}, serta \reg{FLAGS}.
    \item Ulangi langkah 4 dan 5 dengan \file{P} sebagai pembanding,
          lalu bandingkan hasilnya dengan jendela \file{Stack} pada
          Emu8086.
\end{enumerate}

\begin{infobox}
Perintah \file{T} (\eng{trace}) menjalankan satu instruksi dan
\emph{masuk} ke rutin penangan, sedangkan \file{P} (\eng{proceed})
menjalankannya sampai selesai tanpa masuk ke penangan. Keduanya
memanfaatkan \flag{TF} seperti dijelaskan pada
Subbab~\ref{sec:konsep-interupsi}.
\end{infobox}

\subsubsection{C. Tugas 7: Sapaan Berwarna}

Tugas 7 menggabungkan kedua contoh program pada bab ini. Buatlah
program baru bernama \file{tugas7.asm} dengan ketentuan berikut.

\begin{enumerate}
    \item Program meminta nama pengguna dengan \code{INT 21h} fungsi
          \hex{0A} dan kapasitas penyangga \dec{30} karakter, lalu
          meminta NIM pada penyangga terpisah.
    \item Program menetapkan mode video \hex{03} dengan
          \code{INT 10h} fungsi \hex{00}, lalu mencetak baris
          \code{Halo, <nama>} dengan fungsi \hex{13} dan atribut
          kuning.
    \item Program mencetak baris \code{NIM: <nim>} dengan atribut
          hijau terang, lalu menutup tampilan dengan garis merah
          seperti pada \file{warna-layar.asm}.
    \item Program diakhiri \code{INT 21h} fungsi \hex{4C} dan memuat
          komentar berbahasa Indonesia pada setiap baris penting.
\end{enumerate}

\subsection{Hasil yang Dikumpulkan}

\begin{enumerate}
    \item Berkas \file{tugas7.asm} beserta hasil perakitan dan
          tangkapan layar keluarannya.
    \item Tangkapan layar keluaran program pada langkah A.3 beserta
          penjelasan penyebabnya.
    \item Tabel pengamatan nilai \reg{SP} dan isi puncak stack sebelum
          serta sesudah \lstasm|INT 21h|, beserta jawaban langkah B.6
          dan B.7.
\end{enumerate}

\begin{asesmenbox}
\textbf{Asesmen praktikum Bab~\ref{chap:interupsi} (bobot 10\% dari
komponen praktikum).}

\begin{tabular}{@{}p{5.9cm}cp{5.0cm}@{}}
    \toprule
    \textbf{Indikator} & \textbf{Bobot} & \textbf{Kriteria} \\
    \midrule
    Program masukan/keluaran dengan \code{INT 21h} berjalan benar &
    20\% & Ketiga fungsi (\hex{0A}, \hex{09}, \hex{02}) dipakai
    dengan argumen yang tepat \\
    \addlinespace
    Tugas 7 selesai sesuai ketentuan & 35\% & Seluruh ketentuan
    terpenuhi dan teks berwarna tampil benar \\
    \addlinespace
    Ketepatan tabel pengamatan stack & 25\% & Nilai \reg{SP} dan isi
    stack tercatat benar pada setiap langkah \\
    \addlinespace
    Penjelasan hasil pengamatan & 10\% & Alasan perubahan \reg{SP}
    dan makna \file{T} serta \file{P} dijelaskan dengan tepat \\
    \addlinespace
    Kerapian kode dan jurnal praktikum & 10\% & Komentar lengkap dan
    jurnal dapat diulang orang lain \\
    \bottomrule
\end{tabular}
\end{asesmenbox}
